% ==============================================================================
% The Grand Daily Tournament — 2026-09-02 Match (Day 3 Advanced)
% Locked template: EB Garamond + Garamond-Math + STKaiti (v4.1)
% High-Capacity & High-Depth Edition: 3 Math Problems + 2 QM Problems
% ==============================================================================
\documentclass[a4paper,11pt]{article}

\usepackage{fontspec}
\usepackage{amsmath}
\usepackage{unicode-math}
\defaultfontfeatures{Ligatures=TeX}
\setmainfont{EBGaramond}[
  Path           = {c:/texlive/2024/texmf-dist/fonts/opentype/public/ebgaramond/},
  Extension      = .otf,
  UprightFont    = *-Regular,
  ItalicFont     = *-Italic,
  BoldFont       = *-SemiBold,
  BoldItalicFont = *-SemiBoldItalic,
  Ligatures      = TeX,
  Numbers        = OldStyle
]
\setsansfont{LibertinusSans}[
  Path        = {c:/texlive/2024/texmf-dist/fonts/opentype/public/libertinus-fonts/},
  Extension   = .otf,
  UprightFont = *-Regular,
  ItalicFont  = *-Italic,
  BoldFont    = *-Bold,
  Scale       = MatchLowercase
]
\setmathfont{Garamond-Math.otf}[
  Path  = {c:/texlive/2024/texmf-dist/fonts/opentype/public/garamond-math/},
  Scale = MatchLowercase
]

\usepackage{xeCJK}
\xeCJKsetup{
  CJKmath      = false,
  PunctStyle   = kaiming,
  CheckSingle  = true,
  AutoFakeBold = true,
  AutoFakeSlant= false
}
\setCJKmainfont{STKaiti}[
  Path        = {C:/Windows/Fonts/},
  UprightFont = STKAITI.TTF,
  BoldFont    = STKAITI.TTF,
  ItalicFont  = STKAITI.TTF,
  Script      = Default
]
\setCJKsansfont{Microsoft YaHei}
\setCJKmonofont{FangSong}

\usepackage[margin=1.5cm,headheight=14pt,headsep=10pt,footskip=18pt]{geometry}
\usepackage{booktabs,enumitem,xcolor,graphicx,tabularx}

\usepackage[breakable,skins,hooks]{tcolorbox}
\usepackage{tikz}
\usetikzlibrary{calc,arrows.meta}
\usepackage{fancyhdr,lastpage}
\usepackage{microtype}

\definecolor{navy}{HTML}{0F172A}
\definecolor{slate}{HTML}{1E293B}
\definecolor{royal}{HTML}{1E40AF}
\definecolor{mist}{HTML}{64748B}
\definecolor{border}{HTML}{E2E8F0}
\definecolor{paper}{HTML}{F8FAFC}
\definecolor{forest}{HTML}{065F46}
\definecolor{amber}{HTML}{D97706}
\definecolor{wine}{HTML}{991B1B}
\definecolor{ice}{HTML}{EFF6FF}

\linespread{1.08}
\setlength{\parindent}{0pt}
\setlength{\parskip}{3pt plus 1pt minus 1pt}
\setlength{\abovedisplayskip}{4pt plus 1pt minus 1pt}
\setlength{\belowdisplayskip}{4pt plus 1pt minus 1pt}

\pagestyle{fancy}\fancyhf{}
\fancyhead[L]{\textcolor{mist}{\footnotesize\sffamily 2027 Theoretical Physics · 604 Calculus \& 811 Quantum Mechanics}}
\fancyhead[R]{\textcolor{slate}{\footnotesize\sffamily\bfseries Daily Tournament (Intensive) · 2026-09-02}}
\fancyfoot[L]{\textcolor{mist}{\footnotesize\itshape ``Calculate until the very last step.''}}
\fancyfoot[R]{\textcolor{mist}{\footnotesize\sffamily Page \thepage\ of \pageref{LastPage}}}
\renewcommand{\headrulewidth}{0.4pt}
\renewcommand{\headrule}{\color{border}\hrule width\headwidth height\headrulewidth \vskip-\headrulewidth}

\newtcolorbox{briefing}[1]{%
  enhanced, blanker, breakable,
  left=3.5mm, right=2.5mm, top=2.5mm, bottom=2.5mm,
  borderline west={2.5pt}{0pt}{royal},
  colback=paper,
  fonttitle=\sffamily\bfseries\color{slate},
  title={#1},
  attach title to upper={\par\vspace{1.5mm}}
}

\newtcolorbox{problem}[2]{%
  enhanced, breakable,
  colback=white, colframe=border, boxrule=0.5pt, arc=1.5mm,
  left=3.5mm, right=3.5mm, top=2.5mm, bottom=2.5mm,
  fonttitle=\sffamily\bfseries\color{slate},
  title={\raisebox{0pt}{\color{royal}\rule{2.5pt}{9pt}}\hspace{4pt}#1\hfill{\normalfont\footnotesize\color{mist}\itshape #2}},
  colbacktitle=paper, coltitle=slate,
  titlerule=0pt, toptitle=1.5mm, bottomtitle=1.5mm
}

\newtcolorbox{solution}[1]{%
  enhanced, breakable,
  colback=white, colframe=border, boxrule=0.5pt, arc=1.5mm,
  left=3.5mm, right=3.5mm, top=2.5mm, bottom=2.5mm,
  fonttitle=\sffamily\bfseries\color{forest},
  title={\raisebox{0pt}{\color{forest}\rule{2.5pt}{9pt}}\hspace{4pt}#1},
  colbacktitle=paper, coltitle=forest,
  titlerule=0pt, toptitle=1.5mm, bottomtitle=1.5mm
}

\newtcolorbox{debrief}[1]{%
  enhanced, breakable,
  colback=white, colframe=border, boxrule=0.5pt, arc=1.5mm,
  left=3.5mm, right=3.5mm, top=2.5mm, bottom=2.5mm,
  fonttitle=\sffamily\bfseries\color{slate},
  title={\raisebox{0pt}{\color{amber}\rule{2.5pt}{9pt}}\hspace{4pt}#1},
  colbacktitle=paper, coltitle=slate,
  titlerule=0pt, toptitle=1.5mm, bottomtitle=1.5mm
}

\newtcolorbox{warnbox}{%
  enhanced, blanker,
  left=2.5mm, right=2mm, top=1mm, bottom=1mm,
  borderline west={2pt}{0pt}{wine},
  colback=wine!4, fontupper=\footnotesize
}

\newcommand{\checkbox}{\ensuremath{\mdlgwhtsquare}}
\newcommand{\alerttag}[1]{\textcolor{wine}{\sffamily\bfseries [#1]}}

\begin{document}

% ==============================================================================
% PAGE 1: COVER
% ==============================================================================
\begin{titlepage}
\thispagestyle{empty}
\begin{tikzpicture}[remember picture,overlay]
  \fill[paper] (current page.south west) rectangle (current page.north east);

  \fill[navy] (current page.north west) rectangle ($(current page.north east)+(0,-3.4cm)$);
  \fill[royal] ($(current page.north west)+(0,-3.4cm)$) rectangle ($(current page.north east)+(0,-3.55cm)$);
  \fill[amber] ($(current page.north west)+(0,-3.55cm)$) rectangle ($(current page.north east)+(0,-3.62cm)$);

  \node[anchor=west, text=white] at ($(current page.north west)+(2.2cm,-1.5cm)$) {
    {\fontsize{12}{14}\selectfont\sffamily\bfseries THE GRAND DAILY TOURNAMENT \textperiodcentered\ 2027}
  };
  \node[anchor=west, text=white!50] at ($(current page.north west)+(2.2cm,-2.3cm)$) {
    {\fontsize{8.5}{11}\selectfont\sffamily UCAS \textperiodcentered\ HIAIS \textperiodcentered\ THEORETICAL PHYSICS ADVANCED TRAINING DOSSIER}
  };

  \begin{scope}[shift={($(current page.center)+(0,1.2cm)$)}]
    \draw[line width=0.5pt, border!80, dashed] (0,0) circle (3.3cm);
    \draw[line width=0.7pt, border] (0,0) circle (2.5cm);
    \draw[line width=0.4pt, royal!30] (0,0) circle (1.6cm);
    \draw[line width=1.2pt, slate!60] (-3.8,-2.2) -- (3.8,2.2);
    \draw[line width=1.2pt, slate!60] (-3.8,2.2) -- (3.8,-2.2);
    \draw[line width=0.7pt, royal!65, -{Stealth[length=2.5mm]}] (-4.2,0) -- (4.2,0)
      node[right, font=\footnotesize\sffamily\color{mist}] {$x,\, p$};
    \draw[line width=0.7pt, royal!65, -{Stealth[length=2.5mm]}] (0,-3.6) -- (0,3.6)
      node[above, font=\footnotesize\sffamily\color{mist}] {$\psi(x),\, \phi(p)$};
    \draw[line width=1.4pt, royal, smooth, samples=120, domain=-3.2:3.2]
      plot ({\x}, {2.2*exp(-abs(\x))});
    \draw[line width=1.1pt, amber!85!black, dashed, smooth, samples=120, domain=-3.14:3.14]
      plot ({\x}, {1.3*cos(\x*1.57 r)});
    \fill[navy] (0, 2.2) circle (3pt);
    \fill[amber] (0, 1.3) circle (2.2pt);
    \fill[royal] (1.2, 0.66) circle (2.2pt);
    \fill[royal] (-1.2, 0.66) circle (2.2pt);
  \end{scope}

  \node[anchor=center] at ($(current page.center)+(0,-3.4cm)$) {
    \begin{minipage}{0.84\textwidth}
      \centering
      {\fontsize{30}{36}\selectfont\bfseries\color{navy} Daily Tournament}\\[4pt]
      {\fontsize{11}{14}\selectfont\sffamily\bfseries\color{royal}%
        GRAND ARENA OF THEORETICAL PHYSICS \textperiodcentered\ INTENSIVE ARENA}\\[10pt]
      {\color{navy}\rule{0.45\textwidth}{1.2pt}}\\[8pt]
      {\fontsize{10}{13}\selectfont\color{slate}\itshape
        ``604 Advanced Trio: Partition Mean Value $\times$ Uniform Continuity $\times$ Taylor Bound\\
        811 Double Strike: $\delta$-Potential Well $\times$ Symmetric Infinite Well Momentum Wave''\\[3pt]
        {\small\color{mist}\upshape 5-Problem Intensive Battle \textperiodcentered\ Closed-book \textperiodcentered\ 100\,min Real-time Challenge}}
    \end{minipage}
  };

  \node[anchor=south] at ($(current page.south)+(0,1.5cm)$) {
    \begin{minipage}{0.86\textwidth}
      \begin{tcolorbox}[
        enhanced, colback=white, colframe=border, boxrule=0.6pt, arc=2mm,
        left=4mm, right=4mm, top=3mm, bottom=3mm
      ]
        \renewcommand{\arraystretch}{1.25}
        \sffamily\small
        \begin{tabularx}{\linewidth}{@{}p{2.6cm}X@{}}
          \textbf{Date} & 2026-09-02 \quad\textbar\quad Week 1 (Day 3 Advanced)\\
          \textbf{Modules} & \textbf{604}: 3 Problems (M1, M2, M3) \quad\textbar\quad \textbf{811}: 2 Problems (Q1, Q2)\\
          \textbf{Error Protocol} & 5-Factor: Knowledge\,\textbar\, Modelling\,\textbar\, Logic\,\textbar\, Calculation\,\textbar\, Time\\
          \textbf{Standards} & Target: 100/100 Full Marks \quad\textbar\quad Time Budget: 90--100 min\\
        \end{tabularx}
      \end{tcolorbox}
    \end{minipage}
  };
\end{tikzpicture}
\end{titlepage}

% ==============================================================================
% PAGE 2: §1 TACTICAL BRIEFING & WARM-UP RETRIEVAL
% ==============================================================================
\clearpage

\begin{center}
  {\LARGE\bfseries\color{navy} Theoretical Physics Training \textperiodcentered\ Daily Tournament}\\[3pt]
  {\small\sffamily\color{mist} 2026-09-02 \quad\textbar\quad Week 1 (Day 3 Advanced) \quad\textbar\quad Phase: Multi-Problem Power Training}\\[-4pt]
  {\color{navy}\rule{\textwidth}{0.8pt}}\\[-11pt]
  {\color{border}\rule{\textwidth}{0.3pt}}
\end{center}

\vspace{-6pt}

\begin{briefing}{\S 1\quad Tactical Briefing \& Warm-Up Retrieval}
\begin{itemize}[leftmargin=1.5em, itemsep=2pt, topsep=1pt]
  \item \textbf{604 高数深度攻坚 (3 题)}：\textbf{M1} 介值定理划分区间 + 双中值倒数和；\textbf{M2} 半无界区间 $[0,+\infty)$ 一致连续性截断证明；\textbf{M3} 极值点二阶 Taylor 展开估计导数下界 $f''(\xi) \ge 8$。
  \item \textbf{811 量力物理攻坚 (2 题)}：\textbf{Q1} $\delta(x)$ 势阱导数跳变匹配、束缚态能级与几率流密度；\textbf{Q2} 对称一维无限深势阱宇称分类、动量表象波函数与测不准关系。
  \item \textbf{考试规则}：闭卷、独立书写、严禁查阅公式；先在下方草稿区完成 5 分钟白纸公式破冰。
\end{itemize}

\vspace{1mm}
\textbf{Warm-Up Retrieval Checklist} (5\,min): without notes, write on scratch paper:\\
\checkbox\ Cantor Theorem: Continuous on $[a,b] \implies$ Uniformly continuous \quad
\checkbox\ $\delta(x)$ derivative jump condition: $\Delta \psi'(0) = -\frac{2m\alpha}{\hbar^2}\psi(0)$ \quad
\checkbox\ Momentum wave function: $\phi(p) = \frac{1}{\sqrt{2\pi\hbar}}\int_{-\infty}^\infty \psi(x) e^{-ipx/\hbar}dx$
\end{briefing}

\vspace{4pt}

\begin{tcolorbox}[
  enhanced, colback=white, colframe=border, boxrule=0.6pt, arc=2mm,
  left=4mm, right=4mm, top=3mm, bottom=3mm,
  fonttitle=\sffamily\bfseries\color{slate},
  title={\raisebox{0pt}{\color{royal}\rule{3pt}{10pt}}\hspace{4pt}Warm-Up Retrieval Workspace \textbar\ 白纸破冰草稿区},
  colbacktitle=paper, coltitle=slate,
  titlerule=0pt, toptitle=2mm, bottomtitle=2mm
]
\textcolor{mist}{\footnotesize\itshape 5 分钟闭卷盲写：默写 Cantor 定理条件、$\delta$ 势导数跃度、动量表象傅里叶变换定义。}
\vspace{10.2cm}
\end{tcolorbox}

% ==============================================================================
% PAGE 3: §2 CORE PROBLEM SET (5 PROBLEMS)
% ==============================================================================
\clearpage

{\Large\sffamily\bfseries\color{navy} \S 2\quad Core Problem Set (604 Math Trio + 811 QM Double)}\hfill{\small\sffamily\color{mist}\itshape Timed: 100\,min}

\vspace{5pt}

\begin{problem}{M1 \textbar\ 604: Partition Intermediate Value \& Harmonic Mean of Derivatives}{Suggested: 15\,min}
设函数 $f(x)$ 在闭区间 $[0, 1]$ 上连续，在开区间 $(0, 1)$ 内可导，且 $f(0) = 0$，$f(1) = 1$。
证明：必存在两个互不相同的点 $\xi, \eta \in (0, 1)$，使得：
\[
  \frac{1}{f'(\xi)} + \frac{1}{f'(\eta)} = 2.
\]
\end{problem}

\vspace{3pt}

\begin{problem}{M2 \textbar\ 604: Uniform Continuity on Semi-Infinite Interval $[0,+\infty)$}{Suggested: 20\,min}
设函数 $f(x)$ 在区间 $[0, +\infty)$ 上连续，且极限 $\lim_{x \to +\infty} f(x) = A$（$A$ 为有限实常数）。
证明：$f(x)$ 在整个半无界区间 $[0, +\infty)$ 上一致连续。
\end{problem}

\vspace{3pt}

\begin{problem}{M3 \textbar\ 604: Minimum Point \& Second-Order Taylor Bound}{Suggested: 15\,min}
设函数 $f(x)$ 在 $[0, 1]$ 上具有连续的二阶导数，且 $f(0) = f(1) = 0$。已知 $\min_{x \in [0,1]} f(x) = -1$。
证明：必存在 $\xi \in (0, 1)$，使得 $f''(\xi) \ge 8$。
\end{problem}

\vspace{3pt}

\begin{problem}{Q1 \textbar\ 811: 1D $\delta$-Potential Well, Bound State \& Probability Current}{Suggested: 25\,min}
设质量为 $m$ 的粒子在单 $\delta$ 势场 $V(x) = -\alpha \delta(x)$ ($\alpha > 0$) 中运动。
\begin{enumerate}[leftmargin=1.5em,itemsep=1pt,topsep=1pt]
  \item 严格推导波函数在原点处的导数跃变匹配条件 $\psi'(0^+) - \psi'(0^-) = -\frac{2m\alpha}{\hbar^2}\psi(0)$；
  \item 求解束缚态（$E < 0$）的能量本征值 $E_0$ 与归一化波函数 $\psi_0(x)$，并证明为何不存在奇宇称束缚态；
  \item 计算该束缚态的几率流密度 $J(x)$，并阐述一维定态实波函数几率流密度恒为零的物理意义。
\end{enumerate}
\end{problem}

\vspace{3pt}

\begin{problem}{Q2 \textbar\ 811: Symmetric Infinite Well, Parity \& Momentum Wave Function}{Suggested: 25\,min}
一维粒子处于对称无限深方势阱中：$V(x) = 0$ ($|x| < a$)，$V(x) = \infty$ ($|x| \ge a$)。
\begin{enumerate}[leftmargin=1.5em,itemsep=1pt,topsep=1pt]
  \item 按宇称分类写出基态 $\psi_1(x)$（偶宇称）与第一激发态 $\psi_2(x)$（奇宇称）的归一化波函数与能级；
  \item 若体系处于基态，求其动量表象波函数 $\phi(p) = \frac{1}{\sqrt{2\pi\hbar}}\int_{-\infty}^\infty \psi_1(x) e^{-ipx/\hbar} dx$，并指出动量几率密度 $|\phi(p)|^2$ 的最概然值；
  \item 阐释不确定度关系 $\Delta x \cdot \Delta p \ge \hbar/2$ 在该基态下的物理表现。
\end{enumerate}
\end{problem}

% ==============================================================================
% PAGE 4: §3 MODEL SOLUTIONS & COMMENTARY
% ==============================================================================
\clearpage

{\Large\sffamily\bfseries\color{navy} \S 3\quad Model Solutions \& Commentary}

\vspace{4pt}

\begin{solution}{M1 \textperiodcentered\ Partition Intermediate Value \& Lagrange}
由 $f(0)=0, f(1)=1$ 及介值定理，$\exists c \in (0,1)$ 使得 $f(c) = 1/2$。\\
在 $[0,c]$ 上应用拉格朗日中值定理：$f(c)-f(0) = f'(\xi) c \implies \frac{1}{f'(\xi)} = 2c$（$\xi \in (0,c)$）。\\
在 $[c,1]$ 上应用拉格朗日中值定理：$f(1)-f(c) = f'(\eta) (1-c) \implies \frac{1}{f'(\eta)} = 2(1-c)$（$\eta \in (c,1)$）。\\
相加得 $\frac{1}{f'(\xi)} + \frac{1}{f'(\eta)} = 2c + 2(1-c) = 2$。因 $\xi < c < \eta$，必有 $\xi \neq \eta$。证毕！
\end{solution}

\vspace{3pt}

\begin{solution}{M2 \textperiodcentered\ Cantor Theorem \& Semi-Infinite Partition}
$\forall \varepsilon > 0$，因 $\lim_{x\to+\infty} f(x)=A$，$\exists X > 0$ 当 $x \ge X$ 时 $|f(x)-A| < \varepsilon/4$。\\
故当 $x_1, x_2 \ge X$ 时，$|f(x_1)-f(x_2)| \le |f(x_1)-A| + |A-f(x_2)| < \varepsilon/2$。\\
在闭区间 $[0, X+1]$ 上，由 Cantor 定理，$f$ 一致连续，故 $\exists \delta_1 > 0$（不妨取 $\delta_1 < 1$），当 $x_1, x_2 \in [0, X+1]$ 且 $|x_1-x_2|<\delta_1$ 时，$|f(x_1)-f(x_2)| < \varepsilon$。\\
取 $\delta = \delta_1$。若 $x_1, x_2 \in [0, +\infty)$ 且 $|x_1-x_2| < \delta$：若两者均在 $[0, X+1]$，由 Cantor 定理成立；若至少一个 $\ge X+1$，则因 $|x_1-x_2|<1$，两者均 $\ge X$，由极限控制成立。全区间一致连续！
\end{solution}

\vspace{3pt}

\begin{solution}{M3 \textperiodcentered\ Minimum Point \& Second-Order Taylor Bound}
设最小值点为 $x_0 \in (0,1)$，则 $f(x_0) = -1$，$f'(x_0) = 0$。在 $x_0$ 处展开到二阶：\\
$f(0) = f(x_0) + f'(x_0)(0-x_0) + \frac{1}{2}f''(\xi_1)x_0^2 \implies 0 = -1 + \frac{1}{2}f''(\xi_1)x_0^2 \implies f''(\xi_1) = \frac{2}{x_0^2}$。\\
$f(1) = f(x_0) + f'(x_0)(1-x_0) + \frac{1}{2}f''(\xi_2)(1-x_0)^2 \implies f''(\xi_2) = \frac{2}{(1-x_0)^2}$。\\
因为 $\min\{x_0, 1-x_0\} \le 1/2$，必有 $\max\{f''(\xi_1), f''(\xi_2)\} \ge \frac{2}{(1/2)^2} = 8$。证毕！
\end{solution}

\vspace{3pt}

\begin{solution}{Q1 \textperiodcentered\ 1D $\delta$-Potential Well Solution}
(1) 薛定谔方程积分取极限得 $\Delta\psi'(0) = -\frac{2m\alpha}{\hbar^2}\psi(0)$。\\
(2) 偶宇称解 $\psi(x) = Ae^{-\kappa|x|}$ 代入跃变式 $-2\kappa A = -\frac{2m\alpha}{\hbar^2}A \implies \kappa = \frac{m\alpha}{\hbar^2}$。能级 $E_0 = -\frac{\hbar^2\kappa^2}{2m} = -\frac{m\alpha^2}{2\hbar^2}$，归一化常数 $A = \sqrt{\kappa}$。若为奇宇称，$\psi(0)=0 \implies A=0$，无非零解。\\
(3) 定态波函数为实函数，$\psi_0^*=\psi_0 \implies J(x) = \frac{\hbar}{2mi}(\psi_0\psi_0' - \psi_0\psi_0') \equiv 0$。驻波几率无净定向传输。
\end{solution}

\vspace{3pt}

\begin{solution}{Q2 \textperiodcentered\ Symmetric Well \& Momentum Representation}
(1) 偶宇称基态：$\psi_1(x) = \frac{1}{\sqrt{a}}\cos\left(\frac{\pi x}{2a}\right)$，$E_1 = \frac{\pi^2\hbar^2}{8ma^2}$。奇宇称第一激发态：$\psi_2(x) = \frac{1}{\sqrt{a}}\sin\left(\frac{\pi x}{a}\right)$，$E_2 = \frac{\pi^2\hbar^2}{2ma^2}$。\\
(2) 基态动量波函数：$\phi(p) = \frac{1}{\sqrt{2\pi\hbar a}}\int_{-a}^a \cos\left(\frac{\pi x}{2a}\right)e^{-ipx/\hbar}dx = \sqrt{\frac{a}{\pi\hbar}}\frac{\pi\cos(pa/\hbar)}{\pi^2/4 - p^2a^2/\hbar^2}$。\\
当 $p \to 0$ 时，$\phi(0) = \frac{4}{\pi}\sqrt{\frac{a}{\pi\hbar}}$ 取极大值，$|\phi(p)|^2$ 最概然动量为 $p = 0$。\\
(3) 粒子被局域在宽度 $2a$ 的势阱内，动量展开为单缝衍射型谱宽 $\Delta p \sim \hbar/a$，满足测不准原理。
\end{solution}

% ==============================================================================
% PAGE 5: §4 DAILY DEBRIEF & REFLECTION
% ==============================================================================
\clearpage

{\Large\sffamily\bfseries\color{navy} \S 4\quad Daily Debrief \& Reflection (5-Problem Match)}

\vspace{4pt}

\begin{debrief}{Post-Match Scorecard (fill in before sleep, 5\,min)}
\renewcommand{\arraystretch}{1.2}
\sffamily\small
\begin{tabularx}{\linewidth}{@{}p{2.4cm}X@{}}
  \textbf{604 Calculus} & Time: \underline{\hspace{1.2cm}} min \quad\textbar\quad M1: \checkbox\ PASS \quad M2: \checkbox\ PASS \quad M3: \checkbox\ PASS\\
  & Error type: \checkbox\ Knowledge gap \checkbox\ Cantor interval split \checkbox\ Taylor expansion \checkbox\ Time\\
  \midrule
  \textbf{811 QM} & Time: \underline{\hspace{1.2cm}} min \quad\textbar\quad Q1: \checkbox\ PASS \quad Q2: \checkbox\ PASS\\
  & Error type: \checkbox\ $\delta$-jump \checkbox\ Parity symmetry \checkbox\ Fourier transform to $\phi(p)$ \checkbox\ Calculation\\
  \midrule
  \textbf{Summary} & \textbf{Total Score:} \underline{\hspace{1.4cm}} / 100 \quad\textbar\quad Total Time: \underline{\hspace{1.4cm}} min\\
  & \textbf{Spaced retrieval in 3 days:} 9月5日盲写半无界一致连续性证明与动量表象积分\\
\end{tabularx}
\end{debrief}

\vspace{6pt}

\begin{tcolorbox}[
  enhanced, colback=white, colframe=border, boxrule=0.6pt, arc=2mm,
  left=4mm, right=4mm, top=3mm, bottom=3.5mm,
  fonttitle=\sffamily\bfseries\color{slate},
  title={\raisebox{0pt}{\color{forest}\rule{3pt}{10pt}}\hspace{4pt}Key Derivation Insights \& Traps \textbar\ 突破口与避坑沉淀},
  colbacktitle=paper, coltitle=slate,
  titlerule=0pt, toptitle=1.5mm, bottomtitle=1.5mm
]
\sffamily\small
\textbf{Core Breakthrough (今日核心突破口与关键一步)}:\\[26pt]
\rule{\linewidth}{0.3pt}\\[10pt]
\textbf{Fatal Trap \& Common Error (致命陷阱与避坑警示)}:\\[26pt]
\rule{\linewidth}{0.3pt}\\[10pt]
\textbf{Day +3 Action Item (3 天后间隔重做提取的具体题号与断点)}:\\[20pt]
\rule{\linewidth}{0.3pt}
\end{tcolorbox}

\end{document}
